% !TeX spellcheck = nb_NO
% Kommentaren ovenfor gir instruks til editoren din om at den skal bruke norsk stavekontroll. Dette fungerer bl.a. i TeXstudio.

\documentclass[a4paper,norsk]{article}
% Norsk orddeling
\usepackage[norsk]{babel}
% AMS-pakkene er nyttige
\usepackage{amsmath,amsfonts,amssymb,amsthm}
% For å definere lister, som f.eks. deloppgaver
\usepackage{enumitem}
% Inneholder mye nyttig, slik som \coloneqq
\usepackage{mathtools}
% Penere kalligrafiske fonter
\usepackage[utf8]{inputenc}
\usepackage[T1]{fontenc}
\usepackage[a4paper, margin=0.5in]{geometry}
\usepackage{mathabx}
\usepackage{amsthm}
\usepackage{listings}
\usepackage{xcolor}
\usepackage{ifthen}
\usepackage{hyperref}

\usepackage[most]{tcolorbox}
\tcbuselibrary{theorems, breakable}
\tcbset{before upper=\par, after upper=\par} % allows itemize and enumerate inside boxes

\newtcbtheorem[]{definitionbox}{Definisjon}{
  colback=blue!3,
  colframe=blue!70!black,
  fonttitle=\bfseries\color{white},
  coltitle=white,
  boxrule=0.6pt,
  arc=2mm,
  enhanced,
  breakable
}{def}

\newtcbtheorem[]{theorembox}{Teorem}{
  colback=green!3,
  colframe=green!50!black,
  fonttitle=\bfseries,
  coltitle=black,
  boxrule=0.6pt,
  before upper=\par\ignorespaces,
  after upper=\par\unskip,
  arc=2mm,
  enhanced,
  breakable
}{thm}

\newtcbtheorem[]{oppgavebox}{Oppgave}{
  colback=gray!5,
  colframe=black!60,
  fonttitle=\bfseries,
  coltitle=black,
  boxrule=0.6pt,
  arc=2mm,
  enhanced,
  breakable
}{oppg}

\newtcolorbox{sectionbox}[1][]{
  colback=white,
  colframe=black!70,
  boxrule=0.7pt,
  arc=3mm,
  top=2mm, bottom=2mm, left=2mm, right=2mm,
  title=#1,
  enhanced,
  breakable
}

\newtcbtheorem[]{corollarybox}{Korollar}{
  colback=purple!5,
  colframe=purple!60!black,
  fonttitle=\bfseries\color{white},
  coltitle=white,
  boxrule=0.6pt,
  arc=2mm,
  enhanced,
  breakable,
  before upper=\par\ignorespaces,
  after upper=\par\unskip
}{cor}

\newtcbtheorem[]{lemmabox}{Lemma}{
  colback=orange!5,
  colframe=orange!70!black,
  fonttitle=\bfseries,
  boxrule=0.6pt,
  arc=2mm,
  enhanced,
  breakable
}{lem}

\newtcbtheorem[]{propositionbox}{Proposisjon}{
  colback=cyan!5,
  colframe=cyan!60!black,
  fonttitle=\bfseries,
  coltitle=black,
  boxrule=0.6pt,
  arc=2mm,
  enhanced,
  breakable
}{prop}

\newtcbtheorem[]{proofbox}{Bevis}{
  colback=white,
  colframe=black,
  fonttitle=\bfseries\color{white},
  coltitle=white,
  colbacktitle=black,
  title style={opacity=0.9},
  boxrule=0.6pt,
  arc=2mm,
  enhanced,
  breakable
}{proof}

\newlist{suboppgave}{enumerate}{1}
\setlist[suboppgave,1]{
  label=\textbf{(\alph*)},
  ref=\theenumi\alph*,
  leftmargin=1.5em,
  itemsep=0.4em,
  topsep=0.3em,
  parsep=0pt,
  partopsep=0pt
}



\setcounter{secnumdepth}{2}  % Sections and subsections numbered

% Python style  
\lstdefinestyle{pythonstyle}{
    language=Python,
    basicstyle=\ttfamily\small,
    keywordstyle=\color{blue},
    commentstyle=\color{green},
    numbers=left,
    frame=single
}

% Bruk bedre symboler for ≥, ≤, epsilon og phi
\renewcommand{\geq}{\geqslant}
\renewcommand{\leq}{\leqslant}
\newcommand{\eps} {\varepsilon}
\renewcommand{\epsilon}{\varepsilon}
\renewcommand{\phi}{\varphi}

% Definer de vanligste mengdene og funksjonene
\newcommand{\R}{\mathbb{R}}
\newcommand{\K}{\mathbb{K}}
\newcommand{\C}{\mathbb{C}}
\newcommand{\N}{\mathbb{N}}
\newcommand{\Z}{\mathbb{Z}}
\newcommand{\Q}{\mathbb{Q}}
% \L er definert som noe annet fra før av, så vi må redefinere \L
\renewcommand{\L}{\mathcal{L}}
% Rommet av polynomer
\newcommand{\Pol}{{\mathcal{P}}}
\DeclareMathOperator{\id}{id}
\DeclareMathOperator{\Image}{im}
\DeclareMathOperator{\Kernel}{ker}
\DeclareMathOperator{\Span}{span}
\DeclareMathOperator{\Rank}{rank}
\DeclareMathOperator{\Diag}{diag}
\DeclareMathOperator{\Proj}{proj}
\newcommand{\basis}{{\mathcal{B}}}
\newcommand{\basiss}{{\mathcal{C}}}
\newcommand{\basisss}{{\mathcal{D}}}

\usepackage[most]{tcolorbox}
\newtcolorbox{algoritme}[1]{
  title=Algoritme~#1,
  colback=white,
  colframe=black,
  boxrule=0.5pt,
  arc=2mm,
  left=3mm,
  right=3mm,
  top=1mm,
  bottom=1mm
}

% Definer oppgavemiljøet
\newenvironment{oppgave}[1]%
	{\trivlist{}\item\relax\textbf{Løsning på #1.}\medspace}%
	{\endtrivlist}
% Deloppgavelister
\newlist{deloppgave}{enumerate}{1}
\setlist[deloppgave,1]{label=(\textbf{\alph*}),align=left}

\title{Oblig 4 \\ MAT1125, høst 2025}
\author{Oliver Ekeberg}
\begin{document}
\maketitle

\tableofcontents

\section*{5.1 Egenpar og egenrom}
\addcontentsline{toc}{section}{5.1 Egenpar og egenrom}

\begin{oppgavebox}{5.1.2}{}
La $ \lambda $  være en egenverdi for T

a) Vis at hvis u er en tilhørende egenvektor og $ \alpha_{}^{} \in \mathbb{K} \neq 0$  så er også $ \alpha_{}^{} u  $ en egenvektor tilhørende $ \lambda  $

\begin{proofbox}{5.1.2}{}
    $Tu = \lambda u$
    Hvis $ \alpha_{}^{} u $ er en tilhørende egenvektor til $ \lambda  $, så er 
    \[
        T( \alpha_{}^{} u) = \lambda \left( \alpha_{}^{} u \right) 
    \]
    Siden T er en lineær operator, så er 
    \begin{align*}
        T \left( \alpha_{}^{} u \right) &= \alpha_{}^{}  T(u) \\
        &= \lambda \left( \alpha_{}^{} u \right) \\
        &= \alpha_{}^{}  \lambda u
    \end{align*}
\end{proofbox}

b) Vis at hvis u og v er egenvektorer for $ \lambda  $ og $ u+v\neq 0 $, så er $ u+v $ også en egenvektor   

\begin{proofbox}{}
    Hvis $ u+v $ er en egenvektor tilhørende $ \lambda  $, så er 
    \begin{align*}
        T \left( u+v \right) &= \lambda \left( u+v \right)  \\
        &= \lambda u + \lambda v
    \end{align*}
    gitt at $ u+v \neq 0 $
    Egenrommet som består av alle egenvektorer slik at. $ Tu = \lambda u $, inkludert lineær kombinasjoner av egenvektorer. I tillegg så er mengden av alle tilhørende egenverdier til en egenvektor, er ikke helt komplett, fordi den ikke inkluderer $ \lambda =0 $. Og 0 kan være en egenverdi for en matrise. For eksempel, hvis det karakteristiske polynomet $ p_{ \lambda }(t) $ har en felles faktor $ \lambda  $ i alle ledd av polynomet.    
\end{proofbox}
\end{oppgavebox}


\begin{oppgavebox}{5.1.3}{}
La T være en operator med egenverdi $ \lambda  $. Forklar hvorfor følgende er ekvivalent. 
\begin{itemize}
    \item $ \lambda  $ er en egenverdi for T 
    \item $ ker \left( T - \lambda id \right) \neq \left\{ 0 \right\}  $ 
    \item Operatoren $ T - \lambda id $ er ikke injektiv 
\end{itemize}

\begin{proofbox}{5.1.3 i) $ \rightarrow $ ii) }{}
Skal først vise at $ i) \rightarrow ii) $ Hvis $ \lambda  $ er en egenverdi for T, så er $ T u = \lambda u $ . Dette er det samme som å si at $ T u = \lambda id u $. Dette kommer av at 

\[
    id(u) = u
\]
Da kan vi skrive om å si at

\begin{align*}
    Tu &= \lambda id u \\
    Tu - \lambda id u &= 0 \\
    \left( T - \lambda id \right) u &= 0 \\
    S(u) = 0
\end{align*}
Som sier at med $ \lambda  $ som egenverdi, så er tilhørende egenvektor $ u \in ker \left( S \right) $. En egenvektor har egenskapen at uansett, så  $ u \neq 0 $. Dermed kan ikke $ ker S = \left\{ 0 \right\}  $ og det er bevist
\end{proofbox}

\begin{proofbox}{5.1.3 i) $ \rightarrow $ iii) }{}
Skal nå vise at $ i) \rightarrow iii) $ Fra Lemma 2.3.3, så vet vi at det er ekvivalent å si at 

\begin{itemize}
    \item T er injektiv
    \item $ Tu = 0 $ kun hvis $ u=0 $
    \item $ ker T = \left\{ 0 \right\}  $   
\end{itemize}

Vi begynner med at vi har $ T u = \lambda  u $, og at $ \left( T - \lambda  id \right) u = 0  $ som sier at for operatoren $ S = T -id \lambda  $, så er $ u \in kerS $. Siden en egenvektor $ u \neq 0 $, så vet vi fra lemma 2.3.3 at $ ker S \neq \left\{ 0 \right\}  $. Fordi $ u \in  ker S $ og $ u \neq 0 $  
\end{proofbox}

\begin{proofbox}{5.1.3 ii) $ \rightarrow $ iii) }{}
Skal nå vise at $ ii) \rightarrow iii) $ Dette er opplagt, fordi en lineær operator T har at $ ker T = \left\{ 0 \right\}  $ hvis og bare hvis $ Tu = 0 $ kun hvis $ u=0 $. Siden dette ikke er en egenvektor, så er   $ ker T \neq \left\{ 0 \right\}  $. bare erstatt T med $ T - \lambda  id $ og du får at $ T - \lambda id $ ikke er injektiv.  
\end{proofbox}

\begin{proofbox}{5.1.3 iii) $ \rightarrow $ i)}{}
$ iii) \rightarrow i) $ Hvis $ T - \lambda id $ er ikke injektiv, så er ikke $ Tu = 0 $ kun hvis $ u=0 $, og dermed er $ \left( T - \lambda id \right)u = 0  $ for en $ u\neq 0 $ og $ \lambda  $ er dermed er $ \left( \lambda , u \right)  $ et egenpar.
\end{proofbox}

\begin{proofbox}{5.1.3 iii) $ \rightarrow $ ii)}{}
$ iii) \rightarrow ii) $. Om $ T - \lambda id $ er ikke injektiv, så er da $ \left( T - \lambda id \right)u= 0 $ for en $ u\neq 0 $. Dermed er $ ker \left( T - \lambda id \right) \neq \left\{ 0 \right\}   $ 
\end{proofbox}

\begin{proofbox}{5.1.3 ii) $ \rightarrow $ i)}{}
$ ii) \rightarrow i) $. Hvis $ ker \left( T - \lambda id \right) \neq 0 $, så er det en $ \left( T - \lambda id \right)u = 0  $ for en $ u\neq 0 $. Dermed kan vi skrive 
\begin{align*}
    \left( T - \lambda id \right)u &= 0 \\
    Tu - \lambda id u &= \\
    Tu &= \lambda id u \\
    Tu = \lambda u
\end{align*}
som sier at $ \lambda  $ er en egenverdi for T med tilhørende egenvektor $ u $ og egenparet $ \left( \lambda , u \right)  $   
\end{proofbox}
\end{oppgavebox}



\begin{definitionbox}{5.1.6}{}
La T være en operator på et vektorrom U og la $ \lambda  $ være en egenverdi for T. Den \textbf{\textit{geometriske multiplisiteten til T }} er tallet $ \gamma_T ( \lambda ) := dim \left( E_T ( \lambda ) \right)  $  
\end{definitionbox}


\begin{oppgavebox}{5.1.7}{}
Forklar hvorfor en egenverdi $ \lambda  $ har nøyaktig $ \gamma_T ( \lambda ) $ antall tilhørende lineært uavhengige egenvektorer.

\begin{proofbox}{5.1.7}{}
La $ T \in \mathcal{L} (U) $ ha egenverdi $ \lambda  $. Det tilhørende egenrommet til $ \lambda  $ er 
\[
    E_T( \lambda ) := ker \left( T - \lambda id \right) = \left\{ u \in  U : Tu = \lambda u \right\} 
\]
En vektor $ u \in  U $ kan være egenvektor til høyst 'en egenverdi. Dette er fordi hvis vi har 
\begin{itemize}
    \item $ Tu = \lambda _1 u $
    \item $ Tu = \lambda _2 u $  
\end{itemize}
Så er 
\[
    T \left( u-u \right) = \left( 
        \lambda _1 - \lambda 2
     \right) u
\]
som sier at $ T0 = \left( \lambda _1 - \lambda _2 \right)u  $. Siden $ u=0 $ ikke kan være en egenvektor, så må $ \lambda _1 - \lambda _2 = 0 \rightarrow \lambda _1 = \lambda _2 $   

Som følger, vil $ E_T(\lambda_1) \cap E_T (\lambda_2) = \left\{ 0 \right\}  $ 
\end{proofbox}

\begin{proofbox}{5.1.7}{}
Fra forrige definisjon, så vil en egenverdi $ \lambda $ ha $ dim \left( E_T ( \lambda ) \right)  $  egenvektorer. Dvs. at for to lineært uavhnegige egenvektorer $ u_1, u_2 $, så er  
\[
    \alpha_{1}^{} u_1 + ... +  \alpha_{n}^{}  u_n = 0
\]
hvis og bare hvis $ \alpha_{1}^{} =...= \alpha_{n}^{}  =0 $. hvis $ u_1, ...,  u_n $ har samme egenverdi, så er 
\[
    T \left( \alpha_{1}^{} u_1 + ... +  \alpha_{n}^{}  u_n   \right) = \lambda \left( \alpha_{1}^{} u_1 + ... +  \alpha_{n}^{}  u_n \right) 
\]
Siden $ u_1, u_2 \in U $, så er fra definisjonen av egenrommet, så er $ u_1,..., u_n \in E_T ( \lambda  ) $. Dermed vil 
\[
    dim \left( E_T ( \lambda ) \right) = dim \left( u_1,..., u_n \right) = n
\]
Så den geometriske multiplisiteten er n
\end{proofbox}
\end{oppgavebox}




\section*{5.2 Egenverdier på endeligdimensjonale rom}
\addcontentsline{toc}{section}{5.2 Egenverdier på endeligdimensjonale rom}


\begin{oppgavebox}{5.2.1}{}

\begin{proofbox}{5.2.1}{}
Anta først at T er en isomorfi, skal så vise at 0 er ikke en egenverdi som følge av dette.

Siden T er en isomorfi, er den bijektiv, og er også surjektiv og injektiv. Dette betyr at $ Tu = 0 $ kun dersom $ u=0$. Eller $ ker T = \left\{ 0 \right\}  $.

En egenverdi av T er slik at $ Tu = \lambda u $. Men siden $ Tu = 0 $ kun hvis $ u=0 $, som ikke er en egenvektor, så kan ikke 0 være en egenverdi. Anta som en kontradiksjon at $ \lambda = 0 $ er en egenverdi. Det vil si at $ Tu = 0  \cdot u=0$. Men da er $ u \in kerT $ samtidig som at $ u\neq 0 $, som motiser at $ kerT = \left\{ 0 \right\}  $. Dette er en motsigelse.\\




Skal nå vise motsatt vei, at hvis 0 ikke er en egenverdi, så må T være en isomorfi, under forutsetning at $ dimU < \infty $.

Anta at $ dimU < \infty $, og at $ \lambda = 0 $ ikke er en egeverdi. Da er $ kerT = \left\{ 0 \right\}  $ , fordi hvis $ u\neq 0 $ og $ Tu=0 $, så ville 0 vært en egenverdi. Siden T er lineær og $ dimU < \infty $, følger det at T er bijektiv og dermed en isomorfi. Dette gjelder fordi $ dim kerT = 0 $, og $ dimU = dimV $ fra dimensjonssatsen.   
\end{proofbox}
\end{oppgavebox}

\begin{propositionbox}{}
    La T være en operator på et endeligdimensjonalt vektorrom U med basis $ \mathcal{B} $. Da er $ \left( \lambda , u \right)  $  et egenpar for T iff. $ \left( \lambda , \left[ u \right]_{ \mathcal{B}}  \right)  $ er et egenpar for matrisen $ \left[ T \right] _{ \mathcal{B}}^{ \mathcal{B}} $
\end{propositionbox}

\begin{oppgavebox}{5.2.2 Proposisjon}{}
\textbf{\textit{Bevis proposisjon 5.2.2}}
\end{oppgavebox}

\begin{definitionbox}{5.2.3}{}
    T operator på endeligdim. vektorrom U. \textbf{\textit{Det Karakteristiske polynomet}} er en funksjon 
    \[
        P_T ( \lambda ) := det \left( T - \lambda id \right) \forall \lambda \in \mathbb{K}
    \]
\end{definitionbox}

\begin{propositionbox}{5.2.4}{}
La T være en operator på et endeligdim vektorrom U over $ \mathbb{K} $. Da er følgende ekvivalent. 

\begin{itemize}
    \item $ \lambda  $ er en egenverdi for T
    \item $ \lambda  $ er en rot av det karakteristiske polynomet til T. Altså $ p_T( \lambda ) =0 $  
\end{itemize}
\end{propositionbox}

\begin{oppgavebox}{5.2.4}{}
\textbf{\textit{Bevis proposisjon }}

\begin{proofbox}{5.2.4 i)}{}
$ i \rightarrow ii) $ Anta først at $ \lambda  $ er en egenverdi for T. Da er $ Tu = \lambda u $. En annen måte å si dette på er hvis $ u \in ker \left( T - \lambda id \right)  $. Hvis $ \lambda  $ er slik, så er $ ker \left( T - \lambda id \right) \neq 0  $ fordi $ u \neq 0 $. Dermed er $ T - \lambda id $ ikke injektiv.\\

Et krav for at en operator er invertibel, skjer hvis den er bijektiv. En operator er bijektiv hvis den er injektiv OG surjektiv. Siden $ T - \lambda id $  \textbf{\textit{ikke}} er injektiv, så er ikke $ T - \lambda id $  invertibel. Et annet særtrekk med en invertibel operator, er hvis $ det S \neq 0 $. Dermed vet vi at for en $ \lambda  $ som tilfredsstiller $ Tu= \lambda u $, så må $ det \left( T - \lambda id \right) =0 $ for at $ \lambda  $ skal være en egenverdi. 
\end{proofbox}

\begin{proofbox}{5.2.4 ii)}{}
$ ii) \rightarrow i) $.  Hvis $ \lambda  $ er en rot av det karakteristiske polynomet, så vet vi at $ det \left(  T - \lambda id \right) = 0 $. Dermed er ikke $ T - \lambda id $ injektiv (fordi den ikke er invertibel). Det vil si at 
\[
    ker \left( T - \lambda id \right) \neq \left\{ 0 \right\} 
\]
Med andre ord, så $ \exists $ en $ u \in ker \left( T - \lambda id \right)  $ hvor $ u \neq 0 $. Med andre ord, så har $ \lambda  $ en egenvektor u. slik at $ \left( \lambda , u \right)  $    er et egenpar
\end{proofbox}
\end{oppgavebox}


\begin{oppgavebox}{5.2.8}{}
\begin{proofbox}{5.2.8}{}
La $ A \in M_2 (\mathbb{R}) $ være en rotasjonsmatrise
\[
    A = \begin{bmatrix}
        \cos\theta&-\sin\theta \\
        \sin\theta& \cos\theta
    \end{bmatrix} 
\]
Da vil egenverdiene være røttene av likningen 
\[
    det \left( A - \lambda id \right) = \left( \cos\theta - \lambda \right) ^2 + \sin^2\theta
\]

\begin{align*}
\left( \cos\theta - \lambda \right) ^2 + \sin^2\theta &= \cos^2\theta - 2\cos\theta \lambda + \lambda^2 + \sin^2 \theta \\
&= \lambda^2 -2\cos\theta +1=0
\end{align*}
Dette bruker vi abc formel for og får at, $ a=1, b=-2\cos\theta, c=1 $

\begin{align*}
    \lambda &= \frac{ 2\cos\theta }{ 2 } \pm \frac{ \sqrt{ 4 \cos^2\theta -4 } }{ 2 }\\
    &= \cos\theta \pm \sqrt{ \cos^2\theta -1 }
\end{align*}
Siden 
\[
    \cos^2\theta + \sin^2\theta = 1 \rightarrow \sin\theta = \sqrt{ 1- \cos^2\theta } = i \sqrt{ \cos^2\theta - 1 }
\]
Så vi får da at 
\[
    \lambda_1 = \cos\theta + i \sin\theta, \lambda_2 = \cos\theta - i \sin\theta
\]
\end{proofbox}
\end{oppgavebox}

\section*{5.3 Diagonalisering}
\addcontentsline{toc}{section}{5.3 Diagonalisering}


\begin{oppgavebox}{5.3.2}{}
La $ D = diag \left(  \alpha_{1}^{}, ..., \alpha_{n}^{}   \right)  $ være en diagonal matrise. Vis at D er inverterbar iff. $ \alpha_{1}^{} , ..., \alpha_{n}^{} \neq 0  $

\begin{proofbox}{}
Anta først at $ \alpha_{1}^{} , ..., \alpha_{n}^{} \neq 0 $, da $ \exists $  en D slik at 
\[
    D = diag \left( \alpha_{1}^{} , ..., \alpha_{n}^{} \right) 
\]
D er inverterbar om $ det D \neq 0 $. For en diagonal matrise D, så er 
\[
    det D = \alpha_{1}^{} \cdot \alpha_{2}^{} \cdot ... \cdot \alpha_{n}^{} 
\]
Siden $ \alpha_{1}^{} , ..., \alpha_{n}^{} \neq 0 $, så er $ det D \neq 0 $ og D er inverterbar.\\


Anta videre at $ \alpha_{k}^{}  = 0 $ for en $ \alpha_{k}^{} \in \left\{ \alpha_{1}^{} , ..., \alpha_{n}^{} \right\}  $. Da vet vi at 
\[
    D e_k = \alpha_{k}^{} e_k = 0 \cdot e_k = 0
\]
\end{proofbox}
\end{oppgavebox}

\begin{definitionbox}{5.3.3}{}
    La U være et n dim vektorrom. En operator $ T \in \mathcal{L}(U) $ er diagonaliserbar om $ \exists $ en basis $ \mathcal{B} = \left( u_1, ..., u_n \right) $ for U slik at $ \left[ T \right] _{ \mathcal{B}}^{ \mathcal{B}} $ er en diagonalmatrise
\end{definitionbox}

\begin{propositionbox}{5.3.4}{}
La T være en operator på et n dim vektorrom. Da er følgende ekvivalent
\begin{enumerate}
    \item T er diagonialiserbar
    \item T har n lineært uavhengige egenvektorer
\end{enumerate}
\end{propositionbox}

\begin{oppgavebox}{5.3.5}{}
La T være en operator på et n dim vektorrom U, og anta at T har n forskjellige egenverdier. Vis at de tilhørende egenvektorene er lineært uavhengige. Konkluder med at T er diagonialiserbar

\begin{proofbox}{5.3.5}{}
Vi har at $ Tu = \lambda u $. Vi har n forksjellige egenverdier. Anta at $ u_k $ har to forskjellige egenverdier, $ \lambda _k', \lambda _k $ . Da må 

\begin{align*}
    Tu_k &= \lambda _k' u_k\\
    Tu_k &= \lambda _k u_k\\
    T \left( u_k - uk \right) &= T(0) \\
    u_k \left( \lambda _k' - \lambda _k \right) &= T(0)\\
\end{align*}
Derfor må $ \lambda _k' = \lambda _k $ som sier at en egenvektor kan ha en distinkt egenverdi. Dermed må n forksjellige egenverdier ha n forskjellige egenrom $ E_T ( \lambda _k) = ker \left( T - \lambda _k id \right)$. Disse egenrommene må være utspent av n forskjellige egenvektorer. Siden egenvektorene utspenner hvert sitt egenrom, så må de være lineært uavhengige. 

Som en konsekvens, så vil 
\[
    u = \alpha_{1}^{} u_1 + ... + \alpha_{n}^{} u_n
\]
være 
\begin{align*}
    Tu &= T \left( \alpha_{1}^{} u_1 + ... + \alpha_{n}^{} u_n \right) \\
    &= \alpha_{1}^{} Tu1 + ... + \alpha_{n}^{} Tun \\
    &= \alpha_{1}^{} \lambda_1 u_1 + ... + \alpha_{n}^{} \lambda _n u_n
\end{align*}
spenne ut hele U, som har dimensjon n. Derfor har T n lineært uavhengige egenvektorer.


\end{proofbox}
\end{oppgavebox}


\begin{definitionbox}{5.3.6}{}
En kvadratisk matrise $ A \in M_n (\mathbb{K}) $ er \textbf{\textit{diagonaliserbar}} dersom $ \exists $ en inverterbar matrise $ R \in M_n (\mathbb{K}) $ og en diagonalmatrise $ D \in M_n (\mathbb{K}) $ slik at 
\[
    A = RDR^{-1}
\]
\end{definitionbox}

\begin{oppgavebox}{5.3.7}{}
Hva er sammenhengen mellom diagonaliserbarhet av en lineær operator og diagonaliserbarhet av en kvadratisk matrise?

HINT: Sammenlign definisjonene for en operator T og den tilhørende matrisen $ \left[ T \right] _{ \mathcal{C}}^{ \mathcal{C}} $ der $ \mathcal{C} $ er standardbasisen i $ \mathbb{K}^n $ 

\begin{proofbox}{5.3.7}{}
Om en $ T \in \mathcal{L}(U) $ er diagonaliserbar, så er det en basis $ \mathcal{B} = \left( u_1, ..., u_n \right)  $ slik at $ \left[ T \right] _{ \mathcal{B}}^{ \mathcal{B}}  $ er en diagonalmatrise. Da er $ \left[ T \right] _{ \mathcal{B}}^{ \mathcal{B}} = \left( \lambda_1, ..., \lambda _n \right)  $. $ \left[ Tu_j \right] _{ \mathcal{B}} = \left[ T \right] _{ \mathcal{B}}^{ \mathcal{B}} \left[ u_j \right] _{ \mathcal{B}} = \left[ T \right] _{ \mathcal{B}}^{ \mathcal{B}} e_j = diag \left( \lambda _1, ..., \lambda _n \right) e_j  = \lambda _j e_j $ 

dermed så er i vårt tilfelle  \textbf{\textit{KOM TILBAKE TIL DENNE oppgaven}}
\end{proofbox}
\end{oppgavebox}


\section*{5.4 Anvendelser av diagonalisering}
\addcontentsline{toc}{section}{5.4 Anvendelser av diagonalisering}

\begin{oppgavebox}{5.4.1}{}
La $ A \in M_n ( \mathbb{C}) $ og la $ x \in \mathbb{C}^n $. Vis at differenslikningen 

\[
    \begin{cases}
        x_{r+1} = A x_r, \text{ for $ r=0, 1, 2, ... $ }\\
        x_0 = x
    \end{cases}
\]
har løsningen $ x_r = A^r x $ for  $ r \in \mathbb{N} $. Vis at om A er diagonaliserbar, for en diagonalmatrise D, kan vi skrive $ x_r = RD^r R^{-1} $ for $ r=0, 1, 2,... $  

\begin{proofbox}{5.3.7}{}
La $ A \in M_n(\mathbb{C}) $ og la $ x \in \mathbb{C}^n $. Vis at differens likningen
\[
    \text{Løsning} = \begin{cases}
        x_{r+1} = Ax_r\\
        x_0 = x
    \end{cases}
\]
har løsningen 
\[
    x_r = A^r x
\]
Dette kan vises med induksjon. 
\begin{itemize}
    \item $ x_1 = Ax $ 
    \item $ x_2 = Ax_1 = A(Ax) = A^2x $
    \item $ x_3 = Ax_2 = A(A^2x) = A^3x $   
\end{itemize}
Da vil for tilfelle k, $ x_k = A^kx $. Og for et tilfelle til, vil 
\[
    x_{k+1} = A^k x * A = A^{k+1}x
\]
Videre for å vise at om $ A = RDR^{-1} $, så er $ A^r= RD^RR{-1} $ fra tidligere i seksjon 5.4, substituerer vi dette inn i likningen for $ x_r $, får vi
\begin{align*}
    x_r &= A^rx\\
    &= RD^rR^{-1}x
\end{align*}
Og dette gjelder for $ r=0, 1, 2,... $ fordi $ A^0=I_n $, og da vil $ x_0 = Ix $ som bare er initial betingelsen 
\end{proofbox}
\end{oppgavebox}

\begin{oppgavebox}{5.4.2}{}

La $ A = \begin{bmatrix}
    1&3\\
    2&0
\end{bmatrix}  $ og la $ x \in \mathbb{R}^2 $ være vilkårlig. Finn løsningen av differens likningen

\[
    \text{Løsning} = \begin{cases}
        x_{r+1} = Ax_r\\
        x_0 = x
    \end{cases}
\]

\begin{proofbox}{5.4.2}{}
La $ A = \begin{bmatrix}
    1&3\\
    2&0
\end{bmatrix}  $ 

\begin{align*}
    det(A-\lambda I) &= det \left(  \begin{bmatrix}
        1-\lambda& 3 \\
        2 &-\lambda
    \end{bmatrix}  \right) \\
    &= \lambda^2 - \lambda - 6 \\
    \left( \lambda -3 \right) \left( \lambda +2 \right) &=0
\end{align*}

Vi har da at $ \lambda_1 = 3 $, og $ \lambda_2 =-2 $  
$ A-3I: $
\[
    \begin{bmatrix}
        -2&3\\
        2&-3
    \end{bmatrix} = \begin{bmatrix}
        -2& 3\\
        0&0
    \end{bmatrix} 
\]Så $ -2x_1 = -3x_2 \rightarrow x_1 = 3/2 x_2 \rightarrow v_1 = \left[ 3, 2 \right] $  
$ A+2I: $
\[
    \begin{bmatrix}
        3&3\\
        2&2
    \end{bmatrix} = \begin{bmatrix}
        1& 1\\
        0&0
    \end{bmatrix} 
\]Så $ 1x_1 = -1x_2 \rightarrow x_1 = -x_2 \rightarrow v_1 = \left[ -1, 1 \right] $  
\[
    R = \begin{bmatrix}
        3&-1\\
        2&1
    \end{bmatrix}, R^{-1} = \frac{ 1 }{ 5 } \begin{bmatrix}
        1&1\\
        -2&3
    \end{bmatrix} 
\]Så diagonaliseringen er da på formen
\[
    A = \begin{bmatrix}
        3&-1\\
        2&1
    \end{bmatrix} \begin{bmatrix}
        3&0\\
        0&-2
    \end{bmatrix} \begin{bmatrix}
        1&1\\
        -2&3
    \end{bmatrix} \frac{ 1 }{ 5 }
\] Dermed gitt en vilkårlig vektor $ x \in \mathbb{R} $, vil $ x_r $  ha løsning 
\[
    x_r = \frac{ 1 }{ 5 } \begin{bmatrix}
        3^{r+1}-2^{r+1} & 3^{r+1} + 3*2^r\\
        2*3^r+(-2)^{r+2} & 2*3^r + 3 (-2)^{r+1}
    \end{bmatrix} x
\]
\end{proofbox}
\end{oppgavebox}

\section*{5.5 Spektralteoremet}
\addcontentsline{toc}{section}{5.5 Spektralteoremet}

\begin{lemmabox}{5.5.1}{}
La U være et endeligdim indreproduktrom over $ \mathbb{K} $, og la $ T \in \mathcal{U} $ være selvadjungert. Da er alle røtter til $ p_T $ reelle. Spesielt har T minst en reell egenverdi.  
\end{lemmabox}

\begin{lemmabox}{5.5.2}{}
la U være et indreproduktrom over $ \mathbb{K}$, og la $ T \in \mathcal{L}(U) $ være en normal operator. Om $ \left( \lambda , u \right)  $ er et egenpar for T, er $ \left( \overline{ \lambda }, u \right)  $ et egenpar for $ T^* $ 
\end{lemmabox}

\begin{oppgavebox}{5.5.2}{}
Bevis Lemma 5.5.2 med teorem 4.8.5

\begin{proofbox}{5.5.2}{}
Fra teorem 4.8.5 vet vi at for en normal operator T på et indreproduktrom U, så er $ T - \alpha_{}^{} id $ normal for enhver $ \alpha_{}^{} \in \mathbb{K} $. Videre er også $ kerT = ker T^* $


Om $ \left( \lambda , u \right)  $  er et egenpar for T, så er $ Tu = \lambda u \rightarrow \left( T - \lambda id \right)u = 0 $, som sier at $ u \in ker \left( T - \lambda id \right)  $. Siden T er normal, er derfor $ T - \lambda id $ normal for enhver $ \lambda \in \mathbb{K} $ (thm 4.8.5) og $ ker \left( T - \lambda id \right) = ker \left( T^* - \lambda id \right)   $ (thm 4.8.5). Dermed er $ u \in ker \left( T^* - \lambda id \right)  $ og 

\begin{align*}
    \left( T^* - \lambda id \right) u &= 0\\
    T^*u - \lambda id u &= 0 \\
    T^*u = \lambda u
\end{align*}

Siden T er normal, og $ T=T^* $, så vet vi fra $ 4.9.2 b)  $ at T er selvadjungert (fordi T er normal, og enhver T som er normal, er selvadjungert), da er 
\begin{align*}
    \lambda \left\lVert u \right\rVert ^2 &= \lambda \langle u, u \rangle_{  } \\
    &= \langle \lambda  u ,u  \rangle_{  }  \\
    &= \langle Tu, u \rangle_{  } \\
    &= \langle u, Tu \rangle_{  } \text{ fordi T er selvadjungert}\\
    &= \langle u, \lambda u \rangle_{  } \\
    &= \overline{ \lambda } \langle u, u \rangle_{  } \\
    &= \overline{ \lambda } \left\lVert u \right\rVert ^2
\end{align*}
Siden $ \left\lVert u \right\rVert \neq 0 $, så kan vi dele på $ \left\lVert u \right\rVert ^2 $  på begge sider, og si at $ \lambda = \overline{ \lambda } $. Dermed vet vi at 
\[
    T^* u = \lambda u \rightarrow T^* u = \overline{ \lambda }u
\]
Og $ \left( \overline{ \lambda }, u \right)  $  er et egenpar for $ T^* $ 
\end{proofbox}
\end{oppgavebox}

\begin{theorembox}{5.5.3 Spektralteoremet over $ \mathbb{C} $ }{}
La U være et endeligdim indreproduktrom over $\mathbb{C} $ av dimensjon n, og la $ T \in \mathcal{L}(U) $ være normal. Da finnes en ortonormal basis $ \mathcal{B} = \left( u_1, ..., u_n \right)  $ for U slik at 

\[
    Tu = \sum_{ k=1 }^{ n } \lambda _k \langle u, u_k \rangle_{  } u_k \forall u \in U
\]
Dette kalles for \textbf{\textit{egendekomposisjonen til U}}
\end{theorembox}

\begin{proofbox}{\textbf{\textit{theorem 5.5.3 Spektralteoremet over $\mathbb{C}$}}}{}
La $ \left( \lambda , u_1 \right)  $ være et egenpar, normalisert slik at $ \left\lVert u \right\rVert=1  $. Definer nå $ U_2 := span \left( u_1 \right)^{\bot}  $. Vi påstår at om $ w \in U_2 $, så er $ Tw \in U_2 $. Av lemma 5.5.2 så er $ \left( \overline{ \lambda }, u_2 \right)  $ et egenpar for $ T^* $. så 

\begin{align*}
    \langle Tw, u_1 \rangle_{  } &= \langle w, T^*u_1 \rangle_{  } \\
    &= \langle w, \overline{ \lambda }u_1 \rangle_{  } \\
    &= \lambda \langle w, u_1 \rangle_{  } =0
\end{align*}
Så da er $ Tw \in U_2 $. Vi kan dermed se på T som en normal operator på det $ n-1 $  dimensjonale rommet $ U_2 $. Velg nå på samme måte en $ u_2 \in U_2 $ til å være en normalisert egenvektor av denne operatoren. Lar så $ U_3 := span \left( u_1, u_2  \right)^{\bot}  $ og itererer på denne måten n ganger. Dette gir en \textbf{\textit{ortonormal}} liste $ \left( u_1, ..., u_n \right)  $ bestående av egenvektorer for T. Siden $ dimU=n $, må listen også være en basis for U.

For å se (5.5.1) skriver vi $ u = \sum_{ k=1 }^{ n } \langle u, u_k \rangle_{  } u_k $ og får

\[
    Tu = T \left( \sum_{ k=1 }^{ n } \langle u, u_k \rangle_{  } u_k \right) = \sum_{ k=1 }^{ n } \langle u, u_k \rangle_{  } Tu_k = \sum_{ k=1 }^{ n } \langle u, u_k \rangle_{  } \lambda _k u_k
\]
\end{proofbox}


\begin{theorembox}{\textbf{\textit{teorem 5.5.4 Spektralteoremet over $\mathbb{R}$}}}{}
La U være et endelimdim indreproduktrom over $ \mathbb{R} $ av dim n. La $ T \in \mathcal{L}(U) $ være selvadjungert. Da finnes en ortonormal basis $ \mathcal{B} = \left( u_1, ..., u_n \right)  $ for U bestående av egenvektorer for T
\end{theorembox}

\begin{proofbox}{5.5.4}{}
Ganske likt som beviset over $ \mathbb{C} $. Fra Lemma 5.5.1 har T minst en reell egenverdi. Bare velg en av røttene til det karakteristiske polynomet med tilhørende egenvektor $ u_1 $ og iterer på samme vis. 
\end{proofbox}

\begin{corollarybox}{5.5.5 Spektralteoremet for matriser}{}
La $ A \in M_n (\mathbb{C}) $ være en normal matrise. Da finnes en \textbf{\textit{unitær}} matrise $ R \in M_n ( \mathbb{C}) $ og en diagonalmatrise $ D \in M_n( \mathbb{C}) $ slik at \[
    A = RDR^*
\]
Kolonnevektorene i R er egenvektorer for A, og diagonalen D inneholder de tilhørende egenverdiene, med hensyn på posisjon.\\

Altså, om $ \left( \lambda _1, u_1 \right)  $ er et egenpar, så er $ Re_1 = u_1 $ og $ De_1 = \lambda _1 $  
\end{corollarybox}

\begin{proofbox}{5.5.5}{}
$ T: \mathbb{C}^n \rightarrow \mathbb{C}^n $ gitt ved $ T(x) := Ax $ er normal. Av spektralteoremet finnes en ortonormal basis av egenvektorer $ \mathcal{B} = \left( u_1, ..., u_n \right)  $. Av teorem 4.3.8 er da 

\[
    \left[ T \right] _{ \mathcal{B}}^{ \mathcal{B}} = \left( \langle Tu_j, u_i \rangle_{  } \right)—{ij} = \left( \lambda _j \langle u_j, u_i \rangle_{  }  \right)_{ij} = \left( \lambda _j, \delta_{ij} \right)_{ij} =: D
\]
diagonalmatrisen med elementer $ \lambda _1, ..., \lambda _n $. Av korollar 4.3.9 er basisskiftematrisen mellom $ \mathcal{B} og \mathcal{C} $ gitt ved $ \left[ id \right] _{ \mathcal{C}}^{ \mathcal{B}} = R := \left( u_1, ..., u_n \right) $, så 
\[
    A = \left[ T \right] _{ \mathcal{C}}^{ \mathcal{C}} = \left[ id \right] _{ \mathcal{C}}^{ \mathcal{B}} \left[ T \right] _{ \mathcal{B}}^{ \mathcal{B}} \left[ id \right] _{ \mathcal{B}}^{ \mathcal{C}} = RDR^{-1} = RDR^*
\]
Vi brukte i siste steg at R er unitær $ R^{-1} = R^* $ 


\end{proofbox}

\begin{oppgavebox}{5.5.6}{}
Forklar hvorfor spektralteoremet for matriser kan omskrives slik:\\

Enhver normal matrise er diagonaliserbar, og egenvektorene kan velges til å være ortonormal

\begin{proofbox}{5.5.6}{}
Fra oppgavebox 4.8.3 fant vi at alle normale matriser er unitære. Dvs. $ A^{-1} = A^* $. Om A er unitær, så er det fra 4.10.2 ekvivalent med å si at kolonnene til A er ortogonale. Det vil si at det er n egenverdier med multiplisitet. Da finnes det en ortogonalmatrise D, og en matrise R slik at $ A = RDR^{-1} $ og A er diagonaliserbar 
\end{proofbox}
\end{oppgavebox}


\begin{corollarybox}{5.5.7}{}
La $ A \in M_n ( \mathbb{R}) $ være en symmetrisk matrise. Da finnes en ortgonal $ R \in M_n( \mathbb{R}) $ og en diagonalmatrise $ D \in M_n( \mathbb{R}) $ slik at 
\[
    A = RDR^T
\]
\end{corollarybox}

\section*{5.6 Singulærverdidekomposisjon}
\addcontentsline{toc}{section}{5.6 Singulærverdidekomposisjon}

SVD er en videre generalisering av spektralteoremet, som sier at alle normale matriser er diagonaliserbare. For en generell avbildning $ T \in \mathcal{L}(U, V) $ hvor U og V er endeligdim indreproduktrom, så kan T ha en SVD, som kan sammenlignes med en diagonalisering


\begin{lemmabox}{5.6.1}{}
    La U og V være endeligdim indreproduktrom, og la $ T \in \mathcal{L}(U, V) $. Operatoren $ T^*T \in \mathcal{L}(U) $
    
    \begin{itemize}
        \item Er selvadjungert $ \left( T = T^* \right)  $ 
        \item Har relle, ikke negative egenverdier (\textbf{\textit{positiv semidefinitt}})
        \item $ kerT = ker T^*T $ og $ imT^* = imT^*T $  
    \end{itemize}
\end{lemmabox}

\begin{proofbox}{lemma 5.6.1}{}
Skal vise i) at $ T^* T $ er selvadjungert. Det skjer om
\[
    T^*T = \left( T^*T \right) ^*
\]
Vi kan skrive om RHS som \begin{align*}
    \left( T^*T \right) ^* &= T^* \left( T^* \right) ^*\\
\end{align*}
Siden $ \left( T^* \right) ^*  = T $, så er
\[
    T^*T = \left( T^*T \right) ^*
\]
og operatoren $ T^*T $ er selvadjungert (altså symmetrisk for reelle matriser, og $ A^* = \overline{A^T}$ når $ A \in M_n(\mathbb{C}) $ )\\


Skal nå vise ii)

Anta først at $ \left( \lambda , u \right)  $ er et egenpar for $ T^*T \in \mathcal{L}(U) $. Og at $ \lambda \in \mathbb{R} $. Da er $ \lambda ^2 \geq 0 $. 

\begin{align*}
    \left\lVert T^*Tu \right\rVert ^2 &= \langle T^*Tu, T^*Tu \rangle_{  } \\
    &= \langle \lambda u, \lambda u  \rangle_{  }  \\
    &= \lambda \overline{\lambda } \langle u, u  \rangle_{  } \\
    &= \lambda ^2 \left\lVert u \right\rVert ^2 \geq 0
\end{align*}
Det er $ \geq 0 $ på grunn av antagelsene vi gjort


Bevis for iii)

Vi påstår at $ ker T = ker T^*T $. Hvor $ Tu = 0 $ er det opplagt, fordi 
\[
    T^*Tu = T^*0 = 0
\]

T er normal fordi $ T^*T = TT^* $, og fra teorem 4.8.5, så er for enhver $ T \in \mathcal{L}(U) $ normal, og da gjelder det at

\[
    kerT = kerT^*
\]

og det er vist at $ T^* 0 = 0 $ fordi $ Tu = 0 $. Motsatt om $ u \in ker T^*T $ er 

\begin{align*}
    \left\lVert Tu \right\rVert ^2 &= \langle Tu, Tu \rangle_{  } 
    &= \langle T^*Tu, u,  \rangle_{  } \\
    &= \langle 0, u,  \rangle_{  } \\
    &=0 
\end{align*}

så $ Tu = 0 $, det vil si at $ u \in ker T $.\\

Vi påstår videre at $ imT^*T = imT^* $. Av proposisjon 4.7.6 (vi) er $ imT^* = \left( kerT \right) ^{\bot} = \left( ker T^* T \right) ^{\bot} $. Siden $ T^*T $ er normal, så er 

\begin{align*}
    \left( ker T^*T \right)^{\bot} &= \left( \operatorname{im} T^*T \right)^{\bot\bot} \\
    &= imT^*T   
\end{align*}
\end{proofbox}


\begin{oppgavebox}{5.6.2}{}
Bevis de samme (eller tilsvarende) egenskapene for operatoren $ TT^* \in \mathcal{L}(V) $
\end{oppgavebox}

\begin{proofbox}{5.6.2}{}
Bevis for i)

Operatoren $ TT^* $ er selvadjungert, om
\[
    TT^* = \left( TT^* \right) ^*
\]
 
\begin{align*}
    \left( TT^* \right) ^* &= \left( T^* \right) ^* T^* \\
    &= TT^*
\end{align*}\\


bevis for ii)
Anta at $ ( \lambda , u )$ er et egenpar for $ TT^* $. Anta også at $ \lambda \in \mathbb{R} $. Da er 

\begin{align*}
    \left\lVert TT^*u \right\rVert ^2 &= \langle TT^*u, TT^*u  \rangle_{  } \\
    &= \langle \lambda u, \lambda u  \rangle_{  } \\
    &= \lambda  \overline{\lambda } \langle u, u \rangle_{  }\\
    &= \lambda ^2 \left\lVert u \right\rVert ^2 
\end{align*}

Bevis for iii)

Vi påstår at $ ker TT^* = kerT^* $. Hvis $ T^*u= 0 $, er
\begin{align*}
    TT^*u &= T0\\
    &= 0
\end{align*}
Motsatt sier vi at om $ u \in ker TT^* $, er 

\begin{align*}
    \left\lVert T^*u \right\rVert ^2 &= \langle T^*u, T^*u \rangle_{  }\\
    &= \langle u, TT^*u \rangle_{  }  \\
    &= \langle u, 0 \rangle_{  } \\
    &= 0
\end{align*}
Med andre ord, er $ \left\lVert T^*u \right\rVert =0 $, som betyr at $ u \in ker T^* $. Dermed er $ kerTT^* = kerT^* $


Vi påstår så at $ imTT^* = imT $ 

\begin{align*}
    \left( imT \right) ^{\bot} &= kerT^*\\
    &= kerTT^*
\end{align*}
Siden $ TT^*  $ er normal, så er 
\[
    kerTT^* = \left( imTT^* \right) ^{\bot}
\]
og da er $$ \left( imT \right) ^{\bot} = \left( imTT^* \right)^{\bot}  $$. Da er det bare å komplekskonjugere begge sider, og vi får at $ imT = imTT^* $ 
\end{proofbox}



Ifølge tidligere, er $ T^*T $ selvadjungert (symmetrisk) og dermed også normal. Så spektralteoremet garanterer at den har en ortonormal liste av egenvektorer $ \left( u_1, ..., u_n \right)  $ med tilhørende egenverdier $ \left( \lambda _1, ..., \lambda _n \right)  $. Da er 

\begin{align*}
    \left\lVert Tu_j \right\rVert ^2_V &= \langle Tu_j, Tu_j \rangle_{ U } \\
    &= \langle u_j, T^*Tu_j \rangle_{ U } \\
    &= \lambda _j \langle u_j, u_j \rangle_{ U } \\
    &= \lambda _j
\end{align*}
Og vi ser dermed at $ \sqrt{ \left\lVert Tu_j \right\rVert ^2 } = \sqrt{ \lambda _j } $ forteller oss hvor mye T strekker eller krymper hvert element i den ortonormale listen med egenvektorer.

\begin{definitionbox}{5.6.3}{}
La U og V være endeligdimensjonale indreproduktrom over $ \mathbb{K} $, og la $ T \in \mathcal{L}(U, V) $. \textbf{\textit{Singulæreverdiene til T}} er kvadratrøttene 
\[
    \sigma_j := \sqrt{ \lambda _j } (j=1, ..., n
\]
til egenverdiene til $ T^*T $. Multiplisiteten til en singulærverdi er $ \mu_{T^*T}(\lambda _j) $  til $ \lambda _j $


Listen av alle singulærverdier er $ \left( \sigma_1, ..., \sigma_n \right)  $ der $ n := dimU $ ordnet i synkende rekkefølge, slik at $ \sigma_1 \geq \sigma_2 \geq ... \geq \sigma_n $ der hver singulær verdi er gjentatt $ \mu_{T^*T}(\lambda _j )$ ganger 
\end{definitionbox}


\begin{definitionbox}{5.6.4}{}
La $ A \in M_{m \times n} (\mathbb{K}) $. \textbf{\textit{Singulærverdiene til A}}  er kvadratrøttene 

\[
    \sigma_j = \sqrt{ \lambda _j } \left( j=1, ..., n \right) 
\]
til egenverdiene til $ A^*A $ \\

En ekvivalent definisjon er at singulæreverdiene til A er singulærverdiene til den lineære avbildningen $ T: \mathbb{K}^m \rightarrow \mathbb{K}^n $ gitt ved $ Tx := Ax\Sigma  $ 
\end{definitionbox}

\begin{oppgavebox}{5.6.6}{}
Vi har at $ A = \begin{bmatrix}
    4&11&14\\
    8&7&-2
\end{bmatrix}  $ og at 

\begin{align*}
    AA^T &= \begin{bmatrix}
        16 + 11^2 + 14^2 & 32 + 77 - 28\\
        32 + 77 - 28 & 64 + 49 + 4
    \end{bmatrix} \\
    &= \begin{bmatrix}
        333 & 81\\
        81& 117
    \end{bmatrix} 
\end{align*}

Egenverdiene er
\begin{align*}
    det \left( AA^T - \lambda I \right) &= \left( 333- \lambda  \right) \left( 117 - \lambda  \right) \\
    &= \lambda^2 - 117 \lambda - 333 \lambda  + 32400 \\&=
    &= \left( \lambda - 360 \right) \left( \lambda  - 90 \right) 
\end{align*}

Og vi ser åpenbart at løsningen til likningen $ det \left( AA^T - \lambda I \right) = 0  $ er
\begin{itemize}
    \item $ lambda_1 = 360  \rightarrow \sigma_1 = \sqrt{ 360 }$
    \item $ \lambda _2 = 90 \rightarrow \sigma_2 = \sqrt{ 90 } $  
\end{itemize} 
\end{oppgavebox}


\begin{propositionbox}{5.6.7}{}
    La U og V være endeligdim indreproduktrom, og la $ T \in \mathcal{L}(U, V) $
    \begin{itemize}
        \item T er injektiv $\iff 0$ \textbf{\textit{ikke}} er en singulærverdi
        \item Antall positive singulærverdier er lik $ dim \left( imT \right)  $ 
        \item Antall singulærverdier som er lik 0 er lik $ dim \left( kerT \right)  $ 
        \item T er surjektiv iff atnall singulærverdier er lik dim V
    \end{itemize} 
\end{propositionbox}

\begin{proofbox}{bevis for 5.6.7 iii)}{}
Bevis for iii)
Siden antall singulærverdier er lik $ n := dimU $, så er av dimensjonssatsen
\[
    n = dim \left( imT \right) + dim \left( kerT \right) 
\]
Hvis $ dim \left( imT \right) = \# \left\{ j: \sigma_k \neq 0 \right\}  $, så er
\[
    n - \# \left\{ j: \sigma_j \neq 0 \right\} = n - dim \left( imT \right) = dim \left( ker T \right) 
\]
\end{proofbox}

\begin{proofbox}{5.6.7 i)}{}
Av Lemma 2.3.3 er T injektiv $ \iff  kerT = \left\{ 0 \right\} $ og av lemma 5.6.1 iii) er $ kerT = kerT^*T = \left\{ 0 \right\} $ $ \iff $ 0 ikke er en egenverdi av $ T^*T $ , som er ekvivalent med at 0 ikke er en singulærverdi av T
\end{proofbox}

\begin{proofbox}{5.6.7 ii)}{}
Av spektralteoremet er $ dim \left( imT^*T \right)  $ lik antall ikke-null egenverdier av $ T^*T $. Av lemma 5.6.1 iii) er $ dim \left( imT^*T \right) = dim \left( imT^*0 \right)  $ som av proposisjon 4.7.6 vii) er lik $ dim \left( imT \right)  $ 
\end{proofbox}

\begin{proofbox}{5.6.7 iv)}{}
Avbildningen T er surjektiv $ \iff $ $ imT = V $  som av proposisjon 1.4.11 er ekvivalent med at $ dim \left( imT \right) = dimV  $ 
\end{proofbox}

\begin{oppgavebox}{5.6.8}{}
La U være et endeligdim indreproduktrom og la $ T \in \mathcal{L}(U) $ være selvadjungert. Forklar hvorfor singulærverdiene til T er lik absolutt verdien av egenverdiene til T 

\begin{proofbox}{5.6.8}{}
Hvis $ T \in \mathcal{L}(U) $ og er selvadjungert, så er $ T= T^* $. Singulærverdiene til T er lik kvadratrøttene av egenverdiene til $ T^*T $. Det vil si at for $ j = 1, ..., n $, så har alle
\[
    \sigma_j = \sqrt{ \lambda _j^2 } = \left| \lambda_j \right| 
\]
\end{proofbox}
\end{oppgavebox}

\begin{theorembox}{5.6.9 Singulærverdidekomposisjonen}{}
La U og V være indreproduktrom med dimensjin hhv. n og m. La $ T \in \mathcal{L}(U, V) $. Da finnes ortonormale basises $ \mathcal{B} = \left( u_1, ..., u_n \right)  $ for U og $ \mathcal{C} = \left( v_1, ..., v_m \right)  $ for V slik at

\[
    \left[ T \right] _{\mathcal{C}}^{ \mathcal{B}} = diag_{m \times n} \left( \sigma_1, ..., \sigma_r \right) 
\]
eller ekvivalent
\[
    Tu = \sum_{ k=1 }^{ r } \sigma_k \langle u, u_k \rangle_{  } v_k \forall u \in U
\]
\end{theorembox}

\begin{proofbox}{5.6.9}{}
Jeg henviser til boken.
\end{proofbox}

\begin{algoritme}{5.6.10}
    Vi oppsummerer beviset for SVD, som bare er en algoritme for å finne en SVD til en vilkårlig operator $ T $ . jeg oppsummerer den her

\begin{enumerate}
    \item Finn $ T^*T $ 
    \item Finn egenverdiene $ \lambda _1, ..., \lambda _n $ tilhørende $ T^*T $ 
    \item Beregn singulærverdiene $ \sigma_j := \sqrt{ \lambda _j } $ 
    \item Finn en ortonormal liste $ \left( u_1, ..., u_n \right)  $ i U bestående av egenvektorer tilhørende $ \lambda _1, ..., \lambda _n $. Gå frem slik:
    \begin{itemize}
        \item Finn egenvektorer $ \tilde{u_1}, ..., \tilde{u_n} $ tilhørende egenverdiene $ \lambda _1, ..., lambda_n $ på vanlig mpte. Hvis en singulærverdi har multiplisitet 2 eller høyere, må du velge de tilhørende egenvektorene til å være ortogonale. Vi oppnår automatisk at for ulike singulærverdier, så er de tilhørende egenvektorene ortonormale
        \item Normaliser egenvektorene
        \[
            u_j := \frac{ 1 }{ \left\lVert u_j \right\rVert }u_j
        \]
        
    \end{itemize}
\end{enumerate}
\end{algoritme}

\begin{corollarybox}{5.6.11}{}
La U, V, T og r være som over

\begin{enumerate}
    \item $ \left( u_{r+1}, ..., u_n \right)  $ er en basis for $ kerT $ 
    \item $ \left( v_1, ..., v_r \right)  $ er en basis for $ imT $ 
    \item T har rang r
    \item $ \left( \sigma_j^2, u_j \right)  $ for $ \left( j=1, ..., n \right)  $ er egenparene til $ T^*T $ 
    \item $ \left( \sigma_j^2, v_j \right)  $ for $ \left( j=1, ..., m \right)  $ er egenparene til $ TT^* $ 
\end{enumerate}
\end{corollarybox}

\begin{proofbox}{5.6.11 i)}{}
Listen $ \mathcal{B} = \left( u_{r+1}, ..., u_n \right)  $ er ortonormal og ligger i $ kerT^*T  = kerT$ (fra prop. 5.6.1 iii)). Av Prop. 5.6.7 iii) er 
\[
    n - r = \# sing. verdier som er lik 0 = dim \left( kerT \right) 
\]
Så hele $ \mathcal{B}  $ utspenner hele $ kerT $ 
\end{proofbox}

\begin{proofbox}{5.6.11 ii)}{}
Listen $ \mathcal{C} = \left( v_1, ..., v_r \right)  $ er ortonormal og ligger i $ imT $, fordi $ v_j = \frac{ 1 }{ \sigma_j } Tu_j $. Av prop. 5.6.7 ii) er da 
\[
    r = \# positive sing. verdier = dim \left( im T \right) 
\]
Så $ \mathcal{C} $ utspenner hele $ imT $ 
\end{proofbox}

\begin{proofbox}{5.6.11 iii)}{}
Vi vet 
\[
    rankT + dim \left( ker T \right) = dim V
\]

Som kommer fra dimensjonssatsen. Dette er fordi vi vet at $ rank T = dim \left( im T \right)  $ 

Siden vektorene $ \left( v_1, ..., v_r \right)  $ danner en basis for $ imT $, så er $ dim \left( imT \right)  = r \rightarrow rankT = r$. 
\end{proofbox}

\begin{proofbox}{5.6.11 iv)}{}
Hvis $ \left( \sigma_j^2, u_j \right)  $ er egenparene til $ T^*T $ for $ j = 1, ..., n $, så er 
\[
    Tu_j = \sigma_j^2 u_j = \left( \sqrt{ \lambda _j } \right)^2 u_j = \lambda _j u_j 
\]
\end{proofbox}

\begin{proofbox}{5.6.11 v)}{}
For $ j=1, ..., m $ er 
\begin{align*}
    TT^* v_j &= \frac{ 1 }{ \sigma_j } TT^*Tu_j \\
    &= \frac{ 1 }{ \sigma_j } T \left( \sigma _{ j} ^2 u_j \right) \\
    &= \sigma _{ j} ^2 u_j
\end{align*}
så $ \left( \sigma _{ j} ^2, u_j \right)  $ er et egenpar til $ TT^* $ for $ j=1, ..., r $. For $ j= r+1, ..., m $ er $ \sigma _{ j} = 0 $ og de tilhørende vektorene $ v_j $ ligger i $ \left( imTT^* \right)^{\bot}  $. Siden $ TT^* $ er normal, er $ kerTT^* = \left( imTT^* \right) ^{\bot} $ så vi finner at $ v_r+1, ..., v_m \in ker TT^* $, og derfor at 
\[
    TT^* v_j = 0 = \sigma _{ j} v_j \forall j = r+1, ..., m
\]
\end{proofbox}


\begin{oppgavebox}{5.6.12}{}
La $ Tu = \sum_{ k=1 }^{ r } \sigma _{ k} \langle u, u_k \rangle_{  } v_k  $ være en SVD av $ T \in \mathcal{L }(U, V) $ . Vis at $ T^* $ har SVD 
\[
    T^* v = \sum_{ k=1 }^{ r } \sigma _{ k} \langle v, v_k,  \rangle_{  } u_k \forall v \in V
\]


\begin{proofbox}{5.6.12}{}
La $ T \in \mathcal{L}(U, V) $ hvor U og V er definert som i teorem 5.6.9. Da har $ TT^* $ ifølge korollar 5.6.11 egenpar $ \left( \sigma _{ j}^2, v_j  \right)  $  for $ j = 1, ..., m $ og

\[
    TT^*v_j = \sigma _{ j} ^2 v_j
\]
Definer nå $ u_j := \frac{ 1 }{ \sigma _{ j}} T^*v_j $  for $ j= 1, ... r $, hvor $ r \leq m $ slik at $ \sigma _{ 1} , ... \sigma _{ r} > 0 $ og $ \sigma _{ r+1} , ... \sigma _{ m} =0 $.

Da er 
\begin{align*}
    \langle u_j, u_k \rangle_{  } &= \frac{ 1 }{ \sigma _{ j} \sigma _{ k}  } \langle T^*v_j, T^*v_k \rangle_{  } \\
    &= \frac{ 1 }{ \sigma _{ j} \sigma _{ k} } \langle TT^* v_j, v_k \rangle_{  } \\
    &= \frac{ \sigma _{ k} ^2  }{ \sigma _{ k} \sigma _{ j}  } \langle v_k, v_j \rangle_{  } = \delta_{jk}
\end{align*}

For en vilkårlig $ v \in V $ er da

\begin{align*}
    T^* v_j &= T^*  \left(  \sum_{ k=1 }^{ m } \langle v, v_k \rangle_{  } v_k  \right) \\
    &= \sum_{ k=1 }^{ m } \langle v, v_k \rangle_{  } T^*v_k \\
    &= \sum_{ k=1 }^{ r } \langle v, v_k \rangle_{  } \sigma _{ k} u_k + \sum_{ k=r+1 }^{ m } \langle v, v_k \rangle_{  } \sigma _{ k} T^*v_k  \\
    &= \sum_{ k=1 }^{ r } \sigma _{ k} \langle v, v_k \rangle_{  } u_k
\end{align*}

Jeg setter bare en vilkårlig $ v \in V $ som $ v = \sum_{ k=1 }^{ m } \langle v, v_k \rangle_{  } v_k $ 
\end{proofbox}
\end{oppgavebox}

\begin{oppgavebox}{5.6.17}{}
Vis at om $ A \in M_{m \times n}(\mathbb{K}) $ er en diagonalmatrise, er singulær verdiene til A lik absolutt verdien til diagonalelementene til A. 
\begin{proofbox}{5.6.17}{}
Om $ A = diag_{m \times n} \left( \sigma _{ 1} , ..., \sigma _{ n}  \right)  $. Da er 

\[
    A =
    \begin{bmatrix}
    \sigma _{ 1}  & 0      & \cdots & 0      \\
    0      & \sigma _{ 2}  & \cdots & 0      \\
    \vdots &        & \ddots & \vdots        \\
    0      & 0      & \cdots & \sigma _{ n}\\
    \vdots &        &        & \vdots \\
    0      & \hdots &        & 0
    \end{bmatrix}_{m \times n}
\]
Da er $ A^* A $ 
\[
    A^* A =
    A =
    \begin{bmatrix}
    \sigma _{ 1}  & 0      & \cdots & 0      & 0 \\
    0      & \sigma _{ 2}  & \cdots & 0      & 0 \\
    \vdots &        & \ddots &        & \vdots \\
    0      & 0      & \cdots & \sigma _{ n}  & 0
    \end{bmatrix}_{m \times n}
    \begin{bmatrix}
    \sigma _{ 1}  & 0      & \cdots & 0      \\
    0      & \sigma _{ 2}  & \cdots & 0      \\
    \vdots &        & \ddots & \vdots        \\
    0      & 0      & \cdots & \sigma _{ n}\\
    \vdots &        &        & \vdots \\
    0      & \hdots &        & 0
    \end{bmatrix}_{m \times n} = \begin{bmatrix}
        \sigma _{ 1} ^2 & \hdots & 0\\
        \vdots & \ddots & \vdots\\
        0 & \hdots & \sigma _{ n} ^2
    \end{bmatrix} 
\]
Matrisen $ A^*A $ har egenverdier $ \sigma _{ 1}^2, ...,   \sigma _{ n} ^2 $ fordi 
\[
    A^* A e_j = \sigma _{ j}^2 e_j
\]
Og følgende har den singulære verdier 
\[
    \sigma _{ j} := \sqrt{ \sigma _{ j} ^2 }= \left| \sigma _{ j}  \right| 
\]
\end{proofbox}
\end{oppgavebox}

\begin{definitionbox}{5.6.13}{}
Singulærverdiene til en matrise $ A \in M_{m \times n}(\mathbb{K})  $ er kvadratrøttene til egenverdiene til $ A^*A $ 
\end{definitionbox}


\begin{theorembox}{5.6.14}{}
For enhver matrise $ A \in M_{m \times n} (\mathbb{K}) $ finnes unitære matriser $ C \in M_m(\mathbb{K}) $ og $ B \in M_n(\mathbb{K}) $  slik at
\[
    A = CDB^*
\]
der $ D = \operatorname{diag}_{m \times n}\left( \sigma_1, \dots, \sigma_n \right) $. 
De første $r$ kolonnene i $C$ er en basis for $\operatorname{im} A$, 
og de siste $n - r$ kolonnene i $B$ er en basis for $\ker A$.
\end{theorembox}


\begin{algoritme}{5.6.15}
For å beregne SVD av A, lar du T være den tilhørende avbildningen $ T (x) := Ax $, finner SVD av $ Tx = \sum_{ j=1 }^{ r } \sigma _{ j} \left( x \cdot u_j \right) v_j  $  av T og lar
\[
    C = \left( v_1, ..., v_m \right), D = diag_{m \times n} \left( \sigma _{ 1} , ..., \sigma _{ n}  \right), B = \left( u_1, ..., u_m \right) 
\]
\end{algoritme}

\section*{5.7 Anvendelser av SVD}
\addcontentsline{toc}{section}{5.7 Anvendelser av SVD}

\begin{oppgavebox}{5.7.5}{}
La $ A = CDB^* $ være en SVD av A, der A er inverterbar. Vis at SVD en til $ A^{-1} $ er

\[
    A^{-1} = BD^{-1} C^*
\]
der $ D^{-1} = diag \left( 1 / \sigma _{ 1} , ..., 1/ \sigma _{ n}  \right)  $ 
\begin{proofbox}{5.7.5}{}
Fra definisjonen av SVD til matriser, så $ \exists $ unitære $ C \in M_n(\mathbb{K}) $ og $ D \in M_n(\mathbb{K}) $  slik at $ A = CDB^* $. Hvis en matrise er unitær, så er

\[
    A^* = A^{-1}
\]
Dette kommer fra 4.10. Dermed så er 

\begin{align*}
    A^{-1} &= \left( CDB^* \right) ^{-1} \\
    &= \left( B^* \right) ^{-1} D^{-1} C^{-1} \\
    &= \left( B^* \right) ^* D^{-1} C^* \\
    &= B D^{-1} C^*
\end{align*}
Der $ D = D^{-1} = diag \left( 1 / \sigma _{ 1} , ..., 1/ \sigma _{ n}  \right) $ 
\end{proofbox}
\end{oppgavebox}



\end{document}